\documentclass{article}
\usepackage{graphicx} % Required for inserting images
\usepackage[margin=1in]{geometry}

\usepackage{amsmath}
\usepackage{amssymb}
\usepackage{amsfonts}
\usepackage{graphicx}
\usepackage{float}
\usepackage{subcaption}
\usepackage{algorithm}
\usepackage{algorithmic}
\usepackage[
    colorlinks=true,
    linkcolor=black,
    urlcolor=blue,
    citecolor=blue
]{hyperref}

\title{COP4934 Senior Design 1 Scrum \& Meetings }
\date{}

\begin{document}

\maketitle
\tableofcontents
\newpage

\section{Introduction}
\subsection{Sprint Planning}
Sprint Planning is an agile event where the scrum team defines the work to be completed during the upcoming sprint. The team establishes a clear sprint goal, selects high-priority items from the product backlog, and breaks them into technical tasks to form the Sprint Backlog. This meeting aligns the team on what can be delivered by the end of the sprint and how the work will be accomplished.
\begin{itemize}
    \item \textbf{Focus:} Planning the upcoming sprint and defining the sprint goal
    \item \textbf{Goal:} Decide what work will be done in the sprint and how it will be completed.
    \item \textbf{Participants:} Scrum Team (Developers, Product Owner, Scrum Master)
    \item \textbf{Outcome:} Sprint Goal, Sprint Backlog, and a shared understanding of the planned work
\end{itemize}
\subsection{Sprint Review}
Sprint Review is held at the end of a sprint to evaluate the completed product increment. The scrum team demonstrates finished work, gathers feedback from stakeholders, and discusses progress toward the overall product goals. The outcome of the review is an updated product backlog and a shared understanding of what was accomplished and what should be prioritized next.
\begin{itemize}
    \item \textbf{Focus:} The Product Increment (what was built)
    \item \textbf{Goal:} Inspect the increment, get stakeholder feedback, and adapt the Product Backlog
    \item \textbf{Participants:} Scrum Team (Developers, PO, SM) + Stakeholders.
    \item \textbf{Outcome:} Updated Product Backlog, shared understanding of progress, feedback for future work
\end{itemize}

\subsection{Sprint Retrospective}
Sprint Retrospective is an internal scrum team meeting focused on improving the team’s development process. The team reflects on what went well, what challenges were encountered, and how collaboration, tools, or workflows can be improved. The result is a set of actionable improvements that the team commits to applying in the next sprint.

\begin{itemize}
    \item \textbf{Focus:} The Process (how the team works)
    \item \textbf{Goal:} Inspect the team's people, processes, and tools; identify improvements
    \item \textbf{Participants:} Scrum Team only (internal event)
    \item \textbf{Outcome:} Actionable improvements to the team's process/collaboration for the next sprint
\end{itemize}
\section{Meeting Minutes}
\subsection{2-16-2026}
\textbf{Attendance:} Everyone in the team attended the meeting. 
\\During this meeting, we assigned the tasks for sprint 1. In the span of the next two weeks, each role will have their own assigned tasks dedicated for their roles, but we will work towards the overall objective of having a basic skeleton or backbone (low implementation) of our project set up and aim to have it dockerized and able to run on everyone's computer. Went through the Jira and added all the tasks and assigned them. Everyone will be in charge of updating their progress of each task on Jira. Below are the tasks that were assigned to each role:
\begin{itemize}
    \item Cybersecurity/Honeypot Lead:
    \begin{itemize}
        \item Choose honeypot method: decide on cowrite or custom SSH
        \item Deploy first honeypot logging the user, pass, commands, and IP
        \item Define threat boundaries: what attacker can do, what they cannot do, and what is logged
        \item Export logs as JSON or CSV with one file per session
       \item \textbf{ Hard Deliverable:} A running SSH honeypot that logs real attacker sessions.
    \end{itemize}
    \item Backend Systems Engineer:
    \begin{itemize}
        \item Scaffold backend: Fast API, routes, models, services 
        \item Create API: POST, GET
        \item Parse incoming honeypot logs: have framework to implement once backend is fully setup 
        \item Dockerize backend
        \item \textbf{Hard Deliverable:} Backend service that accepts attack logs via HTTP.
    \end{itemize}
    \item Database Engineer:
    \begin{itemize}
        \item Design schema tables: sessions, events, credentials
        \item Implement SQL init script 
        \item Connect backend to accept inserts and test queries
        \item Write sample analytic queries
        \item \textbf{Hard Deliverable:} Postgres with real tables storing attack data.
    \end{itemize}
    \item Frontend Visualization Engineer: 
    \begin{itemize}
        \item Scaffold React App: layout, navigation
        \item Create Views: sessions table and detail page
        \item Connect backend: fetch and render data
        \item Dockerize frontend
        \item \textbf{Hard Deliverable:} Dashboard that displays real or mock sessions.
    \end{itemize}
    \item ML Threat Analysis:
    \begin{itemize}
        \item Define ML Problem: clustering or anomaly detection
        \item Design features: session duration, command count, password entropy
        \item Write pipeline: load data, compute features, plot distributions
        \item \textbf{Hard Deliverable:} Jupyter notebook that processes attack data.
    \end{itemize}
\end{itemize}

\subsection{2-23-2026}
\textbf{Attendance:} Everyone was present during this meeting. 
\\\textbf{Progress Update:}
\\The team discussed individual progress on assigned roles and tasks. Overall, everyone has made significant progress and activelty working towards completion on the current sprint's tasks. We are all consistently updating the Github and Jira as tasks are completed.
\\The backend is fully dockerized and we confirmed that it is running locally on everyone's machines/ The frontend is currently in progress and will be completed by the end of the week. 
\\The team also reviewed requirements for the Design Proposal and Team Contract document that is due this weekend as well as what each of us needs to fill out. 
\\\textbf{Scheduling Update:}
There is a slight scheduling conflict of our meeting times since one of our members has a class during this time, so we agreed to moving our weekly meetings to Mondays at 1:30pm. 

\subsection{3-2-2026}
\textbf{Attendance:} Everyone was present during this meeting. 
\\\textbf{Progress Update:} Everything for sprint 1 was completed. First honeypot was deployed on Cowrie and it can send logs to the backend. We will research further on how to automate this in the future. This current sprint's tasks will be sent after the instructor check ins this week. We will update the Jira with the new tasks and assign them to their correct roles. We will also need to work on the Early status report that is due this weekend and provide MVP, Jira updates, sprint plannings, and Agile event info.

\subsection{3-9-2026}
\textbf{Attendance:} Everyone was present during this meeting. 
\\\textbf{Progress Update:} In this meeting, we talked about the tasks for the next two sprints since they are connected. The main goal is to move Sentinel Grid to a real cloud honey net with a working dashboard and structured attack data. Zach sent out tasks for everyone for week 4 to 7 and we also updated the Jira board sprints and assigned tasks per role. Yama will start on the template for the Preliminary Design document so we can split up everyone's parts. We will work parallel to efficiently get tasks done. Next week is spring break, so we will just have a basic progress update on implementation and documentation. 


\section{Senior Design 1}
\subsection{Sprint 0 : 2/1/2026 - 2/15/2026}
\subsubsection{Goal}
Our goal for Sprint 0 is to set the foundation of our project and to establish a clear semester plan.
\subsubsection{Sprint Planning}
The focus of Sprint 0 is to plan the foundation required to establish a clear semester plan. Our team talked about our skills, preferred programming languages, experience, and what we would each like to do in order to figure out everyone's roles and responsibilities. This will help us figure out what each member needs to develop and learn this semester and clear up any questions before we start implementation. 

\begin{itemize}
    \item What can we deliver by the end of this sprint? 
    \\By the end of this sprint, we will deliver a draft of our semester plan that clearly defines team roles, responsibilities, learning gaps, and requirements of the project.
    \item How will we build it? 
    \\We will have a meeting to collaboratively develop the semester plan including planning weekly meetings. Responsibilities will be assigned based on individual strengths and learning goals, and the plan will be documented and reviewed to ensure alignment with project objectives.
    \item What must be done vs nice:
    \\\textbf{Must be done:} defining team roles, assigning responsibilities, identifying required skills and knowledge
    \\\textbf{Nice to have:} detailed stretch goals and general requirements

\end{itemize}

\subsubsection{Sprint Review}
The team reviewed the following:

\begin{itemize}
    \item What was completed during the sprint:
    \\Weekly meeting times were agreed upon to take place online every Mondays at 12pm and in person meetings will take place as needed. We defined team roles, assigned responsibilities, each member had an idea of required skills and tools they needed to learn.
    \item What worked well:
    \\Communication about everyone's background and skills helped the team assign roles effectively. Collaboration during planning ensured that everyone felt heard and aligned with the project goals and roles.
    \item Feedback received:
    \\Its important to have certain milestones and deliverables for each sprint. It was also suggested that documentation should be kept updated throughout the semester for each member so we can keep track of what everyone is working on.

    \item Updates made to the product backlog (new items, changes in scope, or priority adjustments):
    \\Added tasks for researching required technologies, setting up the GitHub, and creating documentation repository.

\end{itemize}

\subsubsection{Sprint Retrospective}
The team discussed:

\begin{itemize}
    \item What went well during the sprint: 
    \\The team established a clear direction for the semester and ensured everyone understood their responsibilities. Meetings were productive and stayed focused on deliverables and what each person had to do next.
    \item What challenges or obstacles were encountered:
    \\Aligning everyone’s schedules for the weekly meetings and there was a learning curve when it came to learning how to use and understand Jira

    \item What could be improved in future sprints:
    The team should also break large goals into smaller, more manageable tasks to improve tracking and accountability. Each member is able to create their own sub tasks and other things they are working on and add them to the Jira and update them as they complete them
    \item Action items to implement in the next sprint:
    \\Set up the development environment, work on project requirements, assign tasks for the next spring

\end{itemize}


\subsection{Sprint 1 : 2/15/2026 - 3/1/2026}
\subsubsection{Goal}
Our goal for Sprint 1 is to complete the design proposal and team contract with the focus on developing the skills and knowledge required for each team role while working on each respective assigned tasks.

\subsubsection{Sprint Planning}
During the 2/16/2026 meeting, all team members were present and Sprint 1 tasks were assigned. Each team member received role specific responsibilities to work on building a basic skeleton and low implementation of the project. The team aims to have a functional backbone of the system set up, dockerized, and capable of running on all team members' machines by the end of the sprint. All tasks were added to Jira, and each member is responsible for updating their task progress throughout the sprint.

\begin{itemize}
    \item What can we deliver by the end of this sprint? 
    \begin{itemize}
        \item A running SSH honeypot that logs attacker sessions (credentials, commands, IP)
        \item A dockerized backend service that accepts attack logs
        \item A Postgres database with implemented schema tables (sessions, events, credentials) storing attack data
        \item A dockerized React frontend that displays real or mock session data
        \item A Jupyter notebook that processes attack data for clustering or anomaly detection
        \item A basic, integrated project skeleton that runs locally on all team members' machines
    \end{itemize}
    \item How will we build it? 
    \\Each role will focus on completing its assigned  deliverables

    \item What must be done vs nice:
    \\\textbf{Must be done:}
    \begin{itemize}
        \item Dockerize backend and frontend
        \item Ensure the system skeleton runs locally for everyone 
        \item Complete the design proposal and team contract.
    \end{itemize}
    \textbf{Nice to have:}
    \begin{itemize}
        \item Complete all hard deliverables for each role
        \item Initial integration testing between all components
        \item Basic analytics visualizations beyond required views
    \end{itemize}

\end{itemize}

\subsubsection{Sprint Review}
The team reviewed the following:
\begin{itemize}
    \item What was completed during the sprint:
    \begin{itemize}
        \item Design proposal and team contract was completed 
        \item Backend and data base set up and dockerized 
        \item Front end scaffold completed
        \item ML notebooks completed for data analysis and data preprocessing pipeline 
        \item Docker for backend verified to run successfully across everyone's environments locally 
        \item Deployed first honeypot on cowrie that can be compromised via ssh root login with no proper authentication
    \end{itemize}
    \item What worked well:
    \begin{itemize}
        \item Separate github branches and roles allowed for parallel development across backend, frontend, database, and ml
        \item Weekly meetings for updates
        \item Regular communication and working through issues together on Discord
        \item Jira task tracking for progress and accountability
    \end{itemize}
    \item Feedback received:
    \begin{itemize}
        \item Clarity and good documentation needed for other members to know how to run certain parts and ensure reproducibility 
    \end{itemize}
    \item Updates made to the product backlog (new items, changes in scope, or priority adjustments):
    \begin{itemize}
        \item Built off of this sprint to add the next steps and tasks for the next sprint 
        \item Idea of exploring more possibilities of supervised learning for classifications 
    \end{itemize}
\end{itemize}

\subsubsection{Sprint Retrospective}
The team discussed:

\begin{itemize}
    \item What went well during the sprint:
    \begin{itemize}
        \item Every member successfully worked on their roles and tasks while learning the technologies they will be using
        \item We have established a project skeleton earlier than expected 
        \item Communication remained consistent 
    \end{itemize}
    \item What challenges or obstacles were encountered:
    \begin{itemize}
        \item Some of us had to learn how to use docker for the first time 
        \item Limited availability of realistic labeled attack data cause restriction in experimenting with supervised classification- will need to find another way 
    \end{itemize}

    \item What could be improved in future sprints:
    \begin{itemize}
        \item Aim to complete assignments the day before so we all can review the day of submissions
        \item Improved documentation updates on Github for readmes for different parts of the project structure 
    \end{itemize}
    \item Action items to implement in the next sprint:
    \begin{itemize}
        \item Look into supervised learning for classification 
        \item Look into how to connect all the different parts together into a functioning structure (connect back end to ml, to front end, to data, etc.)
        \item Not specifically for next sprint but will need to look into automatic sending logs from honeypot to backend in the future
    \end{itemize}
\end{itemize}


\subsection{Sprint 2 : 3/1/2026 - 3/15/2026}
\subsubsection{Goal}
Establish the core cloud infrastructure for SentinelGrid so we can move beyond a local Docker only setup. By the end of Sprint 2, the team aims to have the backend prepared for cloud deployment, the database configured for real log storage, at least one honeypot connected to the backend log pipeline, and the frontend prepared for integration with live API data. 
Also changed meeting times to mondays at 1:30
\subsection{Sprint Planning}
Sprint 2 focuses on turning SentinelGrid from a local prototype into the first version of a cloud-connected cybersecurity monitoring platform. The priority is to deploy or prepare the backend and database in the cloud, confirm that honeypot logs can be accepted and stored through the API, and ensure the frontend can begin consuming real backend data. Since this is the first infrastructure sprint of the new phase, the emphasis is on reliability, connectivity, and end-to-end data flow rather than advanced analytics or polished UI features.

\begin{itemize}
    \item What can we deliver by the end of this sprint? 
    \begin{itemize}
        \item A cloud-hosted backend deployment or a fully prepared backend deployment environment on a VPS
        \item A PostgreSQL database instance configured for SentinelGrid and connected
        \item A verified log ingestion pipeline where honeypot events can be sent to the backend and stored in the database
        \item At least one deployed honeypot node producing structured logs
        \item Backend API endpoints tested for basic functionality such as log submission and session retrieval
        \item A frontend/dashboard environment connected to backend APIs or prepared for deployment and live integration
        \item Updated documentation for deployment, infrastructure setup, and ingestion workflow
        \item Preprocessing data pipeline for ML and feature extraction 
    \end{itemize}
    
    \item How will we build it? 
      We will work and build each part of the structures in parallel on our own branches on Github. Those relying on other parts needing to be setup first will conduct research on what they need to do/start with a skeleton baseline 

    \item What must be done vs nice:
    \\\textbf{Must be done:}
     \begin{itemize}
        \item Set up cloud infrastructure for backend deployment
        \item Deploy or finalize deployment-ready FastAPI backend
        \item Set up PostgreSQL database and confirm schema creation
        \item Connect backend to database successfully
        \item Insert incoming logs into database correctly
        \item Deploy at least one 
        \item Fix frontend integration with backend APIs
        \item Document setup steps and team workflow
        \item ML pipelines for data and features
    \end{itemize}
    \textbf{Nice to have:}
     \begin{itemize}
        \item Add automatic log forwarding scripts instead of manual testing
        \item Add input validation and stronger API error handling
        \item Start preprocessing modularization for ML pipeline integration
        \item Prepare early visualizations for attack sessions if enough real logs are available
    \end{itemize}

\end{itemize}

\subsubsection{Sprint Review}
The team reviewed the following:

\begin{itemize}
    \item What was completed during the sprint:
    \begin{itemize}
        \item Set up the cloud VPS environment for deploying SentinelGrid services
        \item Successfully deployed the FastAPI backend and connected it to the PostgreSQL database
        \item Implemented backend API endpoints for receiving and storing honeypot logs
        \item Connected the Cowrie honeypot to the backend and verified log ingestion
        \item Continued development of the React dashboard and integrated live API data
        \item Expanded the data preprocessing pipeline and began extracting session-level features for future analysis
        \item Updated project documentation for deployment, backend setup, and database configuration
    \end{itemize}

    \item What worked well:
    \begin{itemize}
        \item Team members were able to work in parallel on separate components with minimal merge conflicts
        \item Cloud deployment allowed everyone to test against the same backend environment
        \item Regular communication through Discord and weekly meetings helped resolve integration issues quickly
    \end{itemize}

    \item Feedback received:
    \begin{itemize}
        \item Continue improving documentation for deployment and environment setup
        \item Ensure all API endpoints are consistently tested before frontend integration
        \item Begin thinking about how collected attack data will be processed and categorized later down the road 
    \end{itemize}

    \item Updates made to the product backlog (new items, changes in scope, or priority adjustments):
    \begin{itemize}
        \item Added tasks for automatic log forwarding from honeypots
        \item Added backend analytics endpoints for session summaries and attacker statistics
    \end{itemize}

\end{itemize}

\subsubsection{Sprint Retrospective}
The team discussed:

\begin{itemize}
    \item What went well during the sprint:
    \begin{itemize}
        \item The backend, database, and infrastructure were successfully connected into a functioning system
        \item Team collaboration remained consistent and everyone completed their assigned sprint tasks
        \item Early integration testing helped identify any issues between project components
    \end{itemize}

    \item What challenges or obstacles were encountered:
    \begin{itemize}
        \item Integrating multiple independently developed components required additional debugging
        \item Some frontend features had to wait until backend APIs were finalized
    \end{itemize}

    \item What could be improved in future sprints:
    \begin{itemize}
        \item Perform integration testing more frequently throughout development
        \item Define API contracts earlier to reduce frontend and backend dependency issues
    \end{itemize}
\end{itemize}

\subsection{Sprint 3 : 3/15/2026 - 3/29/2026}
\subsubsection{Goal}
Expand SentinelGrid from a cloud-connected prototype into a more complete live cybersecurity monitoring platform. By the end of Sprint 3, the team aims to stabilize the deployed infrastructure, support multiple honeypot nodes, improve structured log storage and analytics queries, complete the main dashboard views, and prepare the ML pipeline for anomaly detection, clustering, and backend integration.

\subsubsection{Sprint Planning}
Sprint 3 focuses on system completion, analytics, and stability. After establishing the cloud deployment and initial log-ingestion pipeline in Sprint 2, this sprint will extend SentinelGrid into a more functional honeynet platform with multiple honeypots, improved database performance, new backend analytics endpoints, and a more complete React dashboard. In parallel, the ML/data pipeline will move from preprocessing and baseline feature extraction toward clustering and anomaly scoring so that attack sessions can later be surfaced as higher-level patterns rather than raw logs alone. ML will also explore the possibility of labeling data for supervised classification 

\begin{itemize}
    \item What can we deliver by the end of this sprint? 
    \begin{itemize}
        \item Two deployed honeypot nodes, including one SSH honeypot and one HTTP honeypot
        \item Automatic log forwarding from honeypots into the backend API
        \item A stable PostgreSQL database storing structured live attack data
        \item Improved database schema and indexes for faster log queries
        \item Analytics endpoints for top attacker IPs, attack rates, commands, and session summaries
        \item A working React dashboard with pages such as Live Attack Feed, Session Explorer, and Attack Analytics
        \item Visualizations for attack frequency, common commands, and top sources of activity
        \item Initial clustering and anomaly detection outputs from honeypot sessions and external datasets
        \item Documentation for infrastructure, API usage, database maintenance, and ML integration
    \end{itemize}
    \item How will we build it? 
     We will continue to work in parallel on our own branches on Github. Those relying on other parts needing to be setup first will conduct research on what they need to do/start with a skeleton baseline 

    \item What must be done vs nice:
    \\\textbf{Must be done:}
    \begin{itemize}
        \item Deploy and maintain both honeypots 
        \item Ensure automatic log forwarding works 
        \item Improve raw log schema and add indexes on important fields such as timestamp, IP, and session ID
        \item Implement backend analytics endpoints
        \item Complete core dashboard pages for live logs, sessions, and analytics
        \item Add charts for attack trends and attacker activity
        \item Run clustering or anomaly detection pipeline on collected session data
        \item Define backend ready output format for anomaly or behavior analysis results
    \end{itemize}
    \textbf{Nice to have:}
    \begin{itemize}
        \item Add automatic database backups and export scripts
        \item Create fake attacker interactions such as fake admin logins or fake filesystem responses
        \item Add auto refresh and improved filtering in the dashboard
        \item Add richer session drill-down features in the UI
        \item Tune anomaly scoring for near real time detection
        \item Integrate early anomaly scores directly into dashboard views
        \item Improve UI responsiveness and polish for demo readiness
        \item Possible labeling subset for ML supervised classification
    \end{itemize}
    

\end{itemize}

\subsubsection{Sprint Review}

The team reviewed the following:

\begin{itemize}
    \item \textbf{What was completed during the sprint:}
    \begin{itemize}
        \item Successfully deployed Cowrie (SSH) honeypots
        \item Implemented automatic log forwarding from each honeypot to the FastAPI backend using REST API endpoints
        \item Improved the PostgreSQL database schema by separating logs into structured tables and adding indexes on timestamp, source IP, session ID, and event type to improve query performance
        \item Expanded the React frontend with the Live Attack Feed, Session Explorer, and Attack Analytics pages 
        \item Built the initial machine learning pipeline for preprocessing session data, extracting behavioral features, performing K-Means clustering, and detecting anomalous sessions using Isolation Forest 
        \item Established a standardized backend output format for machine learning results to simplify future dashboard integration
    \end{itemize}

    \item \textbf{What worked well:}
    \begin{itemize}
        \item Parallel development allowed frontend, backend, infrastructure, and machine learning tasks to progress simultaneously with minimal blocking
        \item The modular FastAPI architecture made it straightforward to add new analytics endpoints without affecting existing functionality
        \item The machine learning feature engineering pipeline successfully generated meaningful behavioral features that separated different attacker activities into distinct clusters with the use of kmeans
    \end{itemize}

    \item \textbf{Feedback received:}
    \begin{itemize}
        \item Improve dashboard filtering options to better navigate large numbers of attack sessions
        \item Continue improving the visual design and responsiveness of dashboard components
        \item Increase automation for infrastructure deployment and database maintenance where possible
    \end{itemize}

    \item \textbf{Updates made to the product backlog (new items, changes in scope, or priority adjustments):}
    \begin{itemize}
        \item Added support for additional honeypot services such as FTP, Redis, and MySQL
        \item Prioritize backend performance optimization for larger log volumes
        \item Plan implementation of heuristic attack labeling and decision rules for adaptive honeypot deployment.
    \end{itemize}

\end{itemize}

\subsubsection{Sprint Retrospective}

The team discussed:

\begin{itemize}
    \item \textbf{What went well during the sprint:}
    \begin{itemize}
        \item Infrastructure became considerably more stable, allowing continuous collection of live attack data throughout the sprint
        \item Team members remained productive by working independently on well-defined components while communicating regularly through GitHub and sprint meetings
        \item Machine learning experiments successfully demonstrated that behavioral clustering and anomaly detection could identify different categories of attacker activity
        \item The dashboard evolved into a functional monitoring platform suitable for demonstrating live cyber attacks and analytics
    \end{itemize}

    \item \textbf{What challenges or obstacles were encountered:}
    \begin{itemize}
        \item Some dashboard queries became slower as the amount of collected data increased before database indexes were optimized
        \item Live cloud deployments required additional monitoring and debugging to ensure reliable log forwarding
    \end{itemize}

    \item \textbf{What could be improved in future sprints:}
    \begin{itemize}
        \item Increase automated testing for backend APIs and frontend components before deployment
        \item Improve deployment automation
        \item Expand feature engineering using additional temporal and behavioral metrics to improve machine learning performance.
        \item Continue refining the dashboard interface to make navigation more intuitive for non technical users
    \end{itemize}

    \item \textbf{Action items to implement in the next sprint:}
    \begin{itemize}
        \item Develop heuristic attacker labeling and decision rules for automated attack classification
        \item Expand support for additional honeypot protocols and larger scale deployments
    \end{itemize}

\end{itemize}

\subsection{Sprint 4 : 4/12/2026 - 4/26/2026}
\subsubsection{Goal}
Integrate all major SentinelGrid components and prepare for the Preliminary Design Review. By the end of Sprint 4, the team aims to demonstrate a fully functioning pipeline from honeypot attack collection through backend processing, database storage, planned machine learning analysis + pipeline, with dashboard visualization.

\subsubsection{Sprint Planning}
Sprint 4 focuses on integration, testing, and documentation. After completing the core infrastructure during the previous sprint, this sprint will ensure all currently implemented project components work reliably. The team will also prepare documentation and required information and materials for the the Preliminary Design Document.
We will perpare for Senior Design 2, leaving us with a good place to start once the new semester starts.

\begin{itemize}
    \item What can we deliver by the end of this sprint? 
    
    \begin{itemize}
    \item Stable dashboard displaying live attack information
    \item Backend integration for ML analysis results
    \item Preliminary Design Documentation
    \item Additional testing and bug fixes across the system
    \end{itemize}
   
    \item How will we build it? 
    \\ Continue parallel development while prioritizing component integration, testing, documentation, and bug fixes. Frequent integration meetings will be held to quickly resolve interface and deployment issues. The Preliminary Design Document will be split up between roles and each member will have their designated sections they will need to complete 

  
    \item What must be done vs nice:
    \\\textbf{Must be done:}
    \begin{itemize}
    \item Complete dashboard integration
    \item Perform integration and deployment testing
    \item Complete Preliminary Design documentation
    \end{itemize}
    
    \\\textbf{Nice to have:}
    \begin{itemize}
    \item Improve dashboard responsiveness
    \item Additional analytics visualizations
    \item Better deployment automation
    \item Additional API documentation
    \end{itemize}
    

\end{itemize}

\subsubsection{Sprint Review}
The team reviewed the following:

\begin{itemize}
    \item What was completed during the sprint:
    \begin{itemize}
    \item Integrated backend APIs with frontend dashboard
    \item Connected ML analysis pipeline with backend processing
    \item Completed Preliminary Design documentation 
    \item Fixed multiple integration bugs discovered during testing
    \end{itemize}

    \item What worked well:
    \begin{itemize}
    \item Integration testing greatly improved overall system reliability
    \item Clear communication minimized merge conflicts
    \end{itemize}
    \item Updates made to the product backlog (new items, changes in scope, or priority adjustments):
    \begin{itemize}
    \item Prioritize multiple honeypot deployment
    \item Expand adaptive honeypot research
    \item Improve dashboard usability
    \end{itemize}

\end{itemize}

\subsubsection{Sprint Retrospective}
The team discussed:

\begin{itemize}
    \item What challenges or obstacles were encountered:
    
    \begin{itemize}
    \item Integration testing uncovered unexpected API inconsistencies
    \item Coordinating updates across multiple components required additional effort
    \end{itemize}
    
    \item What could be improved in future sprints:
    \begin{itemize}
    \item Begin integration testing earlier
    \item Improve API documentation
    \item Expand testing
    \end{itemize}
    \item Action items to implement in the next sprint:
    \begin{itemize}
    \item Begin planning the implementing adaptive honeynet features
    \item Expand supported honeypot services
    \end{itemize}
\end{itemize}

\section{Senior Design 2}
\subsection{Sprint 5: 5/13/2026 - 5/27/2026}
\subsubsection{Goal}
Begin Senior Design II by completing checkin and prepare for our Critical Design Review (CDR). The focus will be refining the overall architecture, improving deployment reliability, and planning the remaining features required for the final demonstration.

\subsubsection{Sprint Planning}

\begin{itemize}
    \item What can we deliver by the end of this sprint? 
    \begin{itemize}
    \item CDR preparation 
    \item Improved deployment workflow
    \item Updated project backlog and implementation schedule
    \end{itemize}
   
    \item How will we build it? 
    Continue improving existing components while incorporating review and peer feedback to prepare for the Critical Design Review.
    \\\textbf{Must be done:}

    \begin{itemize}
    \item Prepare for the CDR
    \item Practice Presentation
    \item Improve documentation and plan the rest of SD2
    \end{itemize}
    
    \textbf{Nice to have:}
    
    \begin{itemize}
    \item Improve dashboard appearance
    \item Additional testing automation
    \item begin deployment of additional honeypots
    \end{itemize}
    
\end{itemize}

\subsubsection{Sprint Review}
The team reviewed the following:

\begin{itemize}
    \item What was completed during the sprint:
    \begin{itemize}
        \item Started the preperation of the Critical Design Review presentation.
        \item Refined the deployment process to simplify backend and database setup.
        \item Reviewed the remaining product backlog and prioritized features for the remainder of the semester.
    \end{itemize}

    \item What worked well:
    \begin{itemize}
        \item Team members collaborated effectively on presentation preparation.
        \item Existing project components were stable enough to demonstrate the current system.
        \item Documentation improvements made it easier to explain the system architecture and development progress.
    \end{itemize}
    \item Feedback received:
    \begin{itemize}
        \item Checkin did not go as planned, but we will use the feedback to get back on track 
    \end{itemize}
\end{itemize}

\subsubsection{Sprint Retrospective}
The team discussed:

\begin{itemize}
    \item What challenges or obstacles were encountered:
    \begin{itemize}
            \item Balancing presentation preparation with continued development
            \item Identifying deployment issues across different development environments
        \end{itemize}

    \item Action items to implement in the next sprint:
    \begin{itemize}
        \item Begin implementing the remaining high priority backlog items
        \item Expand honeypot deployment and integration testing
    \end{itemize}
\end{itemize}


\subsection{Sprint 6 : 5/27/2026 - 6/10/2026}
\subsubsection{Goal}
Expand SentinelGrid into a complete adaptive honeynet by deploying multiple honeypot services and containerizing each deployment with Docker for scalable management.

\subsubsection{Sprint Planning}

\begin{itemize}
    \item What can we deliver by the end of this sprint? 
    \begin{itemize}
    \item SSH, HTTP, FTP, Redis, MySQL and SMTP honeypots
    \item Docker deployment for every honeypot
    \item Centralized logging
    \item Unified backend support
    \end{itemize}
   
    \item How will we build it? 
    Develop and deploy each honeypot independently before integrating them into the shared backend infrastructure.

  
    \item Must be done:
    
    \begin{itemize}
    \item Deploy all supported honeypots
    \item Dockerize deployments
    \item Normalize logs
    \item Backend compatibility
    \end{itemize}
    
    \textbf{Nice to have:}
    
    \begin{itemize}
    \item Automatic deployment scripts
    \item Health monitoring
    \item Container scaling
    \end{itemize}

\end{itemize}

\subsubsection{Sprint Review}
The team reviewed the following:

\begin{itemize}
    \item What was completed during the sprint:
    \begin{itemize}
        \item Successfully deployed SSH (Cowrie), HTTP, FTP and will work on others like Redis, MySQL, and SMTP honeypot services in the next sprint 
        \item Containerized each honeypot using Docker for consistent deployment and easier management
        \item Configured centralized log collection to forward events from all honeypots to the backend.
        \item Extended the backend to support multiple honeypot types with a unified logging interface
    \end{itemize}
    \item What worked well:
       \begin{itemize}
        \item Docker significantly simplified deployment and testing across different honeypot services
        \item The modular backend architecture allowed new honeypots to be integrated with minimal changes
        \item Independent development of each honeypot allowed easier testing of each type
    \end{itemize}
    
    \item Feedback received:
    \begin{itemize}
        \item Add monitoring for container health and service availability
        \item Ensure consistent metadata is captured across all honeypot logs for future analytics
    \end{itemize}
    \item Updates made to the product backlog (new items, changes in scope, or priority adjustments):
    \begin{itemize}
        \item Added automated Docker deployment and configuration scripts
        \item Add container health monitoring and automatic restart capabilities
        \item Prioritized adaptive honeypot redistribution using the machine learning decision engine
        \item Planned dashboard enhancements to visualize activity across multiple honeypot types
    \end{itemize}
\end{itemize}

\subsubsection{Sprint Retrospective}
The team discussed:

\begin{itemize}
    \item What challenges or obstacles were encountered:
    \begin{itemize}
        \item Different honeypots produced varying log formats, requiring additional normalization logic
        \item Some Docker networking and port configuration issues delayed deployment
        \item Maintaining consistent configurations across multiple containers required additional coordination
        \item Took a while to get the honeypots up and running 
    \end{itemize}
    

    \item What could be improved in future sprints:
     \begin{itemize}
        \item Automate deployment and configuration using scripts
        \item Improve monitoring and alerting for container failures
        \item Create standardized deployment templates to simplify adding new honeypot services
    \end{itemize}

    \item Action items to implement in the next sprint:
     \begin{itemize}
        \item Integrate the machine learning pipeline with the newly deployed honeypots
        \item Begin developing adaptive honeypot redistribution based on detected attack behavior
        \item Update the frontend to display activity from all deployed honeypot services
        \item Finish the implementation of the variety of honeypots 
    \end{itemize}
\end{itemize}



\subsection{Sprint 7 : 6/10/2026 - 6/24/2026}
\subsubsection{Goal}
Develop behavioral analysis engine by generating heuristic attacker profiles and implementing decision rules that recommend optimal honeypot deployments based on recent attack behavior.


\subsubsection{Sprint Planning}

\begin{itemize}
    \item What can we deliver by the end of this sprint? 
    \begin{itemize}
    \item Behavioral feature engineering
    \item Heuristic attacker labeling
    \item Decision engine
    \item Recommendation API
    \item Dashboard visualization
    \end{itemize}
    
    \item How will we build it? 
\\

    Continue extending the machine learning pipeline by combining clustering, feature engineering, anomaly detection, heuristic rules, and backend integration.

  
    \item What must be done vs nice:
    \\\textbf{Must be done:}
    \begin{itemize}
    \item Behavioral analysis
    \item Decision rules
    \item Backend APIs
    \item Frontend visualization
    \end{itemize}
    
    \textbf{Nice to have:}
    \begin{itemize}
    \item Automatic deployment recommendations
    \item Confidence scores
    \item Recommendation history
    \end{itemize}

\end{itemize}

\subsubsection{Sprint Review}
The team reviewed the following:

\begin{itemize}
    \item What was completed during the sprint:
    \begin{itemize}
        \item Developed a behavioral analysis pipeline that extracts attacker behavior metrics from normalized session data
        \item Implemented heuristic labeling to classify attacker profiles based on clustering results, anomaly detection, and behavioral features
        \item Built a decision engine that recommends appropriate honeypot deployments using predefined decision rules
        \item Integrated recommendation APIs into the backend for real time access results
    \end{itemize}
    \item What worked well:
     \begin{itemize}
        \item Combining feature engineering with clustering and anomaly detection produced meaningful attacker profiles
        \item Backend APIs integrated smoothly with the frontend, enabling near real time visualization of recommendations on a set interval
    \end{itemize}
    \item Feedback received:
    \begin{itemize}
        \item Include confidence scores to indicate the reliability of heuristic classifications
        \item Consider storing historical recommendations to evaluate changes in attacker behavior over time
    \end{itemize}
    \item Updates made to the product backlog (new items, changes in scope, or priority adjustments):
    \begin{itemize}
        \item Add recommendation history and trend analysis
        \item Added confidence scoring for heuristic attacker profiles
        \item Prioritized automatic honeypot deployment based on decision engine recommendations
    \end{itemize}
\end{itemize}

\subsubsection{Sprint Retrospective}
The team discussed:

\begin{itemize}
    \item What challenges or obstacles were encountered:
    \begin{itemize}
        \item Fine tuning heuristic thresholds required multiple iterations to balance accuracy and generalization.
        \item Integrating outputs from clustering and anomaly detection into a unified recommendation 
        \item Processing larger log datasets increased execution time 
    \end{itemize}

    \item What could be improved in future sprints:
    \begin{itemize}
        \item Optimize feature engineering and model execution for improved performance
        \item Expand behavioral features to support additional honeypot services and attack types
    \end{itemize}
    
    \item Action items to implement in the next sprint:
    \begin{itemize}
        \item Implement automatic honeypot redistribution the decision engine
    \end{itemize}
\end{itemize}



\subsection{Sprint 8 : 6/24/2026 - 7/8/2026}
\subsubsection{Goal}
Integrate the adaptive honeypot recommendation engine into the dashboard while providing users with interactive controls for manually managing honeypot deployments and visualizing live attacks.


\subsubsection{Sprint Planning}

\begin{itemize}
    \item What can we deliver by the end of this sprint? 
    \begin{itemize}
    \item Adaptive dashboard
    \item Live world map
    \item Honeypot distribution controls
    \item Deployment recommendations
    \end{itemize}
   
    \item How will we build it?
    \\Connect backend recommendation APIs directly to the React dashboard while expanding visualization capabilities
  
    \item What must be done vs nice:
    \\\textbf{Must be done:}
    \begin{itemize}
    \item Recommendation dashboard
    \item World attack map
    \item Manual deployment controls
    \item Live updates
    \end{itemize}

    \\\textbf{Nice to have:}
    \begin{itemize}
\item Geographic statistics
\item Notification settings
\item Moving around map and zoom

\end{itemize}

\end{itemize}

\subsubsection{Sprint Review}
The team reviewed the following:

\begin{itemize}
    \item What was completed during the sprint:
    \begin{itemize}
        \item Integrated the adaptive honeypot recommendation engine into the dashboard
        \item Developed a live world map displaying attacker locations, ips, event types, and active attack events 
        \item Implemented manual honeypot deployment controls, allowing users redistribute honeypot services
        \item Connected backend recommendation APIs to the frontend, providing live deployment recommendations based on recent attacker behavior
        \item Added automatic dashboard updates to display current honeypot distribution without requiring manual refreshes.
    \end{itemize}
    \item What worked well:
    \begin{itemize}
        \item Backend APIs integrated smoothly with the dashboard, enabling real time communication between the recommendation engine and frontend
        \item Interactive visualizations made it easier to understand attack activity and deployment decisions.
        \item Manual deployment controls provided users with flexibility while maintaining automated recommendations.
    \end{itemize}
    
    \item Feedback received:
    \begin{itemize}
        \item Improve geographic visualizations by providing additional country and regional statistics.
        \item Allow users to filter attacks by honeypot type, attacker profile, and time range.
    \end{itemize}
    
    \item Updates made to the product backlog (new items, changes in scope, or priority adjustments):
     \begin{itemize}
        \item Added geographic analytics and attack trend reporting.
        \item Planned enhanced filtering and search capabilities within the dashboard
    \end{itemize}
\end{itemize}

\subsubsection{Sprint Retrospective}
The team discussed:

\begin{itemize}
    \item What challenges or obstacles were encountered:
    \begin{itemize}
        \item Synchronizing backend recommendations with manual deployment actions required careful state management.
        \item Displaying geographic attack information accurately depended on external IP geolocation services.
        \item Testing overall deployment
    \end{itemize}

    \item What could be improved in future sprints:
    \begin{itemize}
        \item Expand dashboard customization options and filtering capabilities.
        \item Improve scalability for larger numbers of honeypot deployments and concurrent attacks.
    \end{itemize}
    \item Action items to implement in the next sprint:
     \begin{itemize}
        \item Conduct final system testing and performance evaluation.
        \item Complete project documentation and user guides or readmes for github.
        \item Prepare the final project demonstration and presentation.
    \end{itemize}
\end{itemize}


\subsection{Sprint 9 : 7/8/2026 - 7/22/2026}
\subsubsection{Goal}
Prepare SentinelGrid for final demonstration by completing testing, documentation, performance optimization, and presentation materials.

\subsubsection{Sprint Planning}

\begin{itemize}
    \item What can we deliver by the end of this sprint? 
    \begin{itemize}
    \item Final dashboard
    \item Complete documentation
    \item Demo ready system
    \item Final presentation
    \end{itemize}

   
    \item How will we build it? 
    \\ Focus on polishing existing features, fixing bugs, improving usability, and preparing demonstration materials.

  
    \item What must be done vs nice:
    \\\textbf{Must be done:}
    
    \begin{itemize}
    \item Final bug fixes
    \item Documentation
    \item User guide
    \item Presentation
    \item Demo preparation
    \end{itemize}
    
    \\\textbf{Nice to have:}
    \begin{itemize}
    \item UI polish
    \item Performance optimization
    \item Additional analytics
    \item Stretch features if time permits
    \end{itemize}
    
\end{itemize}

\subsubsection{Sprint Review}
The team reviewed the following:

\begin{itemize}
    \item What was completed during the sprint:
    \begin{itemize}
        \item Completed end to end testing of the SentinelGrid platform, including the frontend, backend, machine learning pipeline, and Dockerized honeypot infrastructure
        \item Fixed remaining bugs and improved system stability, performance, and overall usability
        \item Finalized technical documentation, installation instructions, architecture diagrams, and user guides.
        \item Prepared the final project presentation and demonstration showcasing adaptive honeypot deployment, behavioral analysis, and live attack monitoring
        \item Performed performance evaluations to verify the system is within acceptable execution times
    \end{itemize}
    \item What worked well:
     \begin{itemize}
        \item Comprehensive documentation made deployment and system maintenance more straightforward.
        \item Integration testing confirmed reliable communication between the dashboard, backend services, machine learning engine, and honeypot containers.
    \end{itemize}
\end{itemize}

\subsubsection{Sprint Retrospective}
The team discussed:
\begin{itemize}
    \item What went well during the sprint:
     \begin{itemize}
        \item Successfully delivered all planned project objectives and integrated every major component into a cohesive platform
        \item Final testing validated the stability and reliability of the complete SentinelGrid system 
        \item Documentation and user guides improved the accessibility and maintainability of the project
    \end{itemize}

\end{itemize}



\end{document}
